%% ============================================================
\section{Problem Formulation and System Overview}\label{sec:overview}
%% ============================================================

\subsection{From Target-Directed Invariants to Loop Summaries}
\label{sec:overview:task}

Fix a loop program $P = (\phi_{\mathrm{pre}}, B, S)$ over integer variables $V$: a precondition $\phi_{\mathrm{pre}}$, a guard $B$, and a body $S$.
A \emph{state} $\sigma : V \to \mathbb{Z}$ is a valuation of $V$ (together with the pre-state values $v_0$ of parameters, so that summaries may relate current to initial values).
Let $\mathcal{R}(P)$ denote the set of states reachable \emph{at the loop head} from some initial state satisfying $\phi_{\mathrm{pre}}$.

\begin{definition}[Inductive invariant]
A predicate $I$ over $V$ is an inductive invariant of $P$ if
(i) \emph{initiation}: every loop-head state reached directly from $\phi_{\mathrm{pre}}$ satisfies $I$; and
(ii) \emph{consecution}: $\{I \wedge B\}\, S\, \{I\}$.
Equivalently, $I \supseteq \mathcal{R}(P)$ and $I$ is closed under the loop's transition relation.
\end{definition}

\begin{definition}[Target-directed generation]
Given $(P, Q)$, produce any inductive $I$ with $I \wedge \neg B \Rightarrow Q$.
\end{definition}

\begin{definition}[Inductive loop summary generation (ours)]
Given $P$ alone, produce an inductive $I$ that is as \emph{tight} as practically expressible: among predicates in the assertion language, $I$ should exclude as much of the unreachable space $\overline{\mathcal{R}(P)}$ as possible.
The ideal object is the strongest inductive invariant, i.e.\ (the assertion-language best approximation of) $\mathcal{R}(P)$ itself.
\end{definition}

The summary task subsumes the target-directed one: any $Q$ provable from \emph{some} inductive invariant is provable from the strongest one.
It is also strictly harder to game.
In the target-directed task, $Q$ itself (weakened to pass initiation) is often an acceptable answer; in the summary task there is no $Q$ to parrot, and weak answers---$\mathit{true}$, the guard, loose boxes---are exactly what the reward of $\S$\ref{sec:method} is designed to score low.

In practice a summary is a \emph{set} of ACSL clauses
$\texttt{loop invariant } e_1; \dots; \texttt{loop invariant } e_k;$
whose conjunction is the invariant; typical summaries pair one or two relational laws (e.g.\ $x = y \cdot z$ preserved by the body) with progress bounds (e.g.\ $0 \le i \le n$).
Soundness is never entrusted to the learner: every clause set is reduced to its greatest inductive subset by Houdini pruning~\citep{flanagan2001houdini} over Frama-C/WP~\citep{framac} before it is scored or reported, so the system's output is certified by construction.

\subsection{System Overview}
\label{sec:overview:system}

\sam{} decomposes into three components that share a common core (a lightweight C/ACSL program model and a safe ACSL predicate evaluator) but are independently deployable:

\begin{stepbox}[1. Example sampler (\S\ref{sec:method:sampler})]
Compiles and executes the instrumented loop over a deterministic, precondition-respecting input sweep; emits \emph{positive} loop-head states (real trace states) and \emph{truthful negative} trace units (provably impossible histories) along a relational and a bound dimension.
The sampler sees only $(\phi_{\mathrm{pre}}, B, S)$---never the postcondition.
\end{stepbox}

\begin{stepbox}[2. Reward framework (\S\ref{sec:method:reward})]
Scores a \emph{group} of rollouts (candidate clause sets sampled for the same program, as in group-relative RL~\citep{shao2024deepseekmath}).
Each rollout is reduced to its surviving inductive core, then scored by the fraction of negative traces it rejects (base), its ablation gain over the rest of the group (marginal), and a junk discount; scores are minimized across independent and held-out sampling seeds.
Packaged as a slim containerized HTTP service (\texttt{POST /reward}) so any RL trainer can consume it, plus an offline JSONL/Parquet scorer.
\end{stepbox}

\begin{stepbox}[3. Inference framework (\S\ref{sec:method:inference})]
Prompts the policy \emph{closed-book} (postcondition stripped) for the strongest ACSL invariants; unions $n$ rollouts, deduplicates, Houdini-prunes, re-annotates, and certifies with Frama-C/WP, re-rolling on failure.
Supports cloud APIs and in-process batched vLLM for RL evaluation.
\end{stepbox}

A fourth ingredient---three adversarial evaluation gates (discrimination, hack sweep, mislabel audit)---is not part of the serving path but guards the reward's integrity as a regression suite ($\S$\ref{sec:experiments}).

\paragraph{Running example.}
Consider a benchmark loop (abridged from our linear suite):
\begin{lstlisting}
void foo(int n) {  // requires n >= 0
  int x = 0, y = n;
  while (y > 0) { x = x + 1; y = y - 1; }
}
\end{lstlisting}
The sampler observes loop-head states such as $(x{=}0, y{=}5)$, $(x{=}1, y{=}4)$, \dots for $n{=}5$ and many other inputs.
Relational negatives perturb dense trace neighborhoods---$(x{=}2, y{=}4)$ with $n{=}5$ breaks $x + y = n$; bound negatives continue past the real exit---$(x{=}6, y{=}{-}1)$ preserves $x+y=n$ yet is impossible.
The summary $\{x + y = n,\; 0 \le y \le n\}$ (with $n \ge 0$ from the precondition) rejects every negative trace and is seed-indifferent, scoring $\approx 1$; the guard restatement $y \ge 0$ rejects only bound negatives on one axis; a sprayed disequality farm rejects canonical negatives but collapses to the minimum on the held-out seed.
Any exit-time fact a user later cares about---$x = n$, $y = 0$, $x \ge y$---follows from the summary, though the generator never saw such a target.
